\label{sec:background}

\paragraph{Memory-Based Meta-RL.} The objective of a meta-RL agent is to adapt its policy to maximize returns over a horizon of one or more attempts in a new task. We focus on a subset of ``black-box'' approaches that reduce meta-learning to the problem of training sequence model policies with end-to-end RL \cite{beck2023survey}. We view adaptation as a form of partial observability \cite{ghosh2021generalization}, where instead of inferring the current state $s_t$ from a sequence of observations $(o_0, \dots, o_t)$, we try to infer the state \textit{and} the identity of the task from trajectories ($\tau$) of observations, actions ($a$), and rewards ($r$): $\tau_{0:t} = (o_0, a_0, r_0, o_1, \dots, a_{t-1}, r_{t-1}, o_t)$ \cite{zintgraf2022fast}. A task's identity is an unobservable variable describing aspects of its dynamics that vary across the training distribution \cite{hallak2015contextual}. Because improved estimates of the current task would increase returns, meta-reasoning can arise implicitly in the latent space of a sequence model that is trained by standard RL updates. 

This sequence-based meta-RL framework, which began with RL$^2$ \cite{wang2016learning, duan2016rl, botvinick2019reinforcement}, has two key advantages. First, training policies on trajectory sequences turns standard partial observability and zero-shot generalization \cite{kirk2022survey} into special cases. In practice, this creates flexible agents that blur formal boundaries between memory, generalization, and meta-learning \cite{benjamins2022contextualize}. Second, end-to-end RL enables self-improvement without demonstrations. Many meta-RL algorithms improve performance in applications where task identification could benefit from additional assumptions \cite{humplik2019meta, liu2021decoupling, rakelly2019efficient, beck2024splagger}. Another line of work recovers the flexibility of ``in-context'' sequence learning when a dataset of demonstrations makes self-improvement unnecessary \cite{laskin2022context, lee2024supervised, raparthy2023generalization, shi2024cross}. Most of the algorithmic variety in meta-RL comes from these two areas. In contrast, the RL$^2$ approach has primarily become an engineering problem. Important developments include increasing the adaptation horizon by enabling long-context sequence models in an RL setting \cite{mishra2017simple, melo2022transformers, lu2023structured}, and a shift away from on-policy updates in favor of off-policy actor-critics \cite{fakoor2019meta, yang2021recurrent, ni2022recurrent}. Combinations of these details allow Transformers with long-term recall to self-improve on large datasets of recycled data \cite{team2023human, ni2023transformers, grigsby2024amago}.

\paragraph{Multi-Task Optimization.} Multi-task learning addresses a trade-off where conflicting objectives decrease performance relative to single-task learning despite an increase in training data. One approach uses specialized optimizers that edit gradients to balance tasks' local optima \cite{shaham2023causes, du2018adapting, yu2020gradient, li2021robust, wang2020gradient, liu2021conflict}. However, it can be as effective (and less expensive) to rescale each task's contribution to the overall loss \cite{xin2022current, hessel2019multi, royer2024scalarization, kurin2022defense, sodhani2021multi}. These techniques require the ability to identify the current task, which is commonly available in multi-task settings: we can view MTRL as a meta-RL problem where the ground-truth identity of each environment has been one-hot labeled in advance. MTRL often depends on task labels in more subtle ways, including networks with task-specific output heads \cite{kumar2022offline} and datasets that balance sampling \cite{d2024sharing} or normalize rewards \cite{garage} across tasks. 

\paragraph{Problem Setting.} We focus on hybrid generalization problems that require elements of meta-learning, multi-task learning, and long-term memory. Meta-RL often evaluates performance over $k > 1$ attempts or episodes. A ``$k$-episode'' objective maximizes the \textit{total} return over all $k$ attempts and creates an exploration/exploitation trade-off at test-time \cite{zintgraf2022fast}. A ``$k$-\textit{shot}'' objective allows for an exploration window where only the final episode counts towards the agent's score \cite{liu2021decoupling, stadie2018some}. We measure generality by dividing a set of task variants ($\mathcal{V}$) into train and test sets. Meta-RL benchmarks are (informally) characterized by their relative density --- meaning $|\mathcal{V}|$ is large and composed of subtle changes to a common objective. Multi-task problems are created by training a single policy to maximize $N$ qualitatively distinct objectives, such as learning to play $N$ Atari games  \cite{bellemare2013arcade}. More general settings might be described as containing $N$ adaptation problems with their own set of variants $\mathcal{V}_0, \dots, \mathcal{V}_N$. These multi-task meta-learning problems are created by attempts to expand training sets by combining existing environments.


\paragraph{Scaling Meta-RL Beyond the Multi-Task Barrier.} Early meta-RL experiments focused on adapting to minor variations of multi-armed bandits \cite{wang2016learning}, gridworlds \cite{stadie2018some}, and locomotion benchmarks \cite{rakelly2019efficient, finn2017model, rothfuss2018promp}. Despite significant algorithmic progress, these domains are still representative of recent research \cite{beck2024splagger, laskin2022context, ni2022recurrent, rimon2024mamba}. Meta-RL requires generating many variations of a core problem while keeping those changes partially observed; it is difficult to design a cohesive benchmark that supports the dense level of variety it takes to induce meta-learning \cite{team2023human, wang2021alchemy}. We can create toy meta-RL problems, but there is a wide gap in complexity before we reach naturally occurring domains where long-term adaptation is necessary. For example, one application of meta-RL would be to learn to master unfamiliar video games over several minutes or hours. Building such an agent would likely involve scaling to a large training set of games. It is difficult to make gradual progress towards this goal because it begins as a memory problem, becomes a multi-task problem, and then only requires meta-learning at an extreme scale. Individual games often require long-term memory, but as we add games, there will inevitably be isolated groups with unique objectives that introduce challenges from MTRL. Eventually, we reach a level of training diversity where adaptation becomes critical and meta-learning emerges. The central issue considered by this paper is that multi-task optimization limits train-time performance at task counts far smaller than would be necessary for test-time generalization to new tasks. A promising direction is to take a method that is already capable of both memory and meta-learning and then find a way to break the barrier where progress is limited by multi-task optimization. However, if we want to adapt to many tasks, we should avoid solutions that scale with the number of training tasks or depend on our ability to identify them manually. 

This ``multi-task barrier'' impacts both long-term applications and current research benchmarks. For a more practical example, consider Meta-World \cite{yu2020meta}. Meta-World is a suite of $50$ robotic manipulation tasks. Each task creates a meta-RL problem by adjusting an unobservable goal so that agents must interpret incoming reward signals to adapt to the current objective. Meta-World ML45 extends generalization by merging meta-RL tasks into a multi-task training set $\cup_{i=0}^{N=45} \mathcal{V}_i$. Identifying the current task during training should be a special case of meta-RL because the changes between tasks are less ambiguous than the changes in goal location between variants of the same task. For example, an agent might distinguish a pick-and-place task from a door opening task after a single observation of the objects on the table. ML45 reserves five tasks to benchmark meta-RL's long-term goal of adapting to entirely new tasks at test-time. However, learning task adaptation from just $N=45$ training tasks is not realistic without additional assumptions \cite{yu2020meta, anand2022procedural}. ML45 is clearly bottlenecked by MTRL: a similar setup \textit{with the meta-RL component removed} is a challenging MTRL benchmark (MT50) by itself \cite{yu2020gradient, liu2021conflict}. MTRL methods use task labels to balance loss functions, separate network architectures, or edit task gradients. However, these techniques scale with $N$ and require a formal identification of each ``task'' that meta-RL otherwise should not need. Our goal is to overcome the same challenge while preserving the meta-RL perspective that task differences are arbitrary so that we can scale to unstructured domains where $N$ is large \cite{reed2022a} or infinite. Therefore, even though our experiments will have a clearly defined $N$, we will enforce a strict constraint where the agent can have zero knowledge of which task is active or how many tasks there are ($N$, $|\mathcal{V}_n|$).